% problem-000306 7 氯酸钾热分解催化
\section*{7. 氯酸钾热分解催化}
$\mathrm { K C l O _ { 3 } }$ 热分解是实验室制取氧⽓的⼀种⽅法。 $\mathrm { K C l O _ { 3 } }$ 在不同的条件下热分解结果如下：

\textit{（原卷此处含表格/仪器试剂表，详见 PDF 或 MD 源文件。）}

已知： (1) $\mathrm { K } ( \mathrm { s } ) + 1 / 2 \mathrm { C l } _ { 2 } ( \mathrm { g } ) = \mathrm { K C l } ( \mathrm { s } )$ $\Delta H ^ { \ominus } \ ( 1 ) = \ - 4 3 7 \ \mathrm { k J \ m o l ^ { - 1 } }$ (2) $\mathrm { K ( s ) } + 1 / 2 \mathrm { C l } _ { 2 } + 3 / 2 \mathrm { O } _ { 2 } ( \mathrm { g } ) = \mathrm { K C l } \mathrm { O } _ { 3 } ( \mathrm { s } )$ $\Delta H ^ { \ominus } \ ( 2 ) = \mathrm { - } 3 9 8 \mathrm { k J m o l ^ { - 1 } }$ (3) $\mathrm { K ( s ) + 1 / 2 \ C l _ { 2 } + 2 \ O _ { 2 } ( g ) = K C l O _ { 4 } ( s ) }$ $\Delta H ^ { \ominus }$ (3) = −433 kJ $\mathrm { m o l ^ { - 1 } }$

\subsection*{7-1}
根据以上数据，写出上述三个体系对应的分解过程的热化学⽅程式。

\subsection*{7-2}
⽤写 $\mathrm { M n O _ { 2 } }$ 催化 $\mathrm { K C l O _ { 3 } }$ 分解制得的氧⽓有轻微的刺激性⽓味，推测这种⽓体是什么，并提出确认这种⽓体的实验⽅法。
